\subsection{Hierarchical Constraints for Concept Lattices}
\label{sec:constrained_latticed}

In the following, we prove that \ac{mlb} constraints induce a valid closure operator and that the resulting lattice encapsulates the search space of all valid dendrograms. 

\subsubsection{Theoretical Analysis of MLB Constraints}
% \label{sec:theoretical_analysis}
Let $G$ be the set of documents. A set of \ac{mlb} constraints is given as $\Sigma \subseteq G \times G \times G$. A constraint $MLB_{x y z} = (d_x, d_y, d_z) \in \Sigma$ expresses a hierarchical dependency: whenever documents $d_x$ and $d_z$ occur together in a cluster, $d_y$ must also belong to that cluster.

\subsubsection{MLB Constraints as a Closure Operator}
To utilize \ac{fca}, we must first establish that these constraints define a closure system.

\begin{definition}[$\Sigma$-Consistency]
A subset $A \subseteq G$ is called $\Sigma$-consistent (or valid) if for all $(d_x, d_y, d_z) \in \Sigma$, the condition $\{d_x, d_z\} \subseteq A \implies d_y \in A$ holds.
\end{definition}

\begin{theorem}
The mapping $\phi_{\Sigma}: \mathcal{P}(G) \to \mathcal{P}(G)$, which maps any set $A$ to the smallest $\Sigma$-consistent set containing $A$, is a closure operator.
\end{theorem}

\begin{proof}
Let $\mathcal{S}_{\Sigma}$ be the family of all $\Sigma$-consistent subsets of $G$. The set $G$ itself is trivially $\Sigma$-consistent (as it contains all elements). Furthermore, the intersection of any family of $\Sigma$-consistent sets is also $\Sigma$-consistent:
Let $\{C_i\}_{i \in I}$ be a family of sets where each $C_i \in \mathcal{S}_{\Sigma}$. Let $D = \bigcap_{i \in I} C_i$.
Suppose $\{d_x, d_z\} \subseteq D$ for some constraint $(d_x, d_y, d_z) \in \Sigma$. By definition of intersection, $\{d_x, d_z\} \subseteq C_i$ for all $i \in I$. Since each $C_i$ is $\Sigma$-consistent, $d_y \in C_i$ for all $i \in I$. Consequently, $d_y \in D$. Thus, $D$ is $\Sigma$-consistent.

Since $\mathcal{S}_{\Sigma}$ is closed under arbitrary intersection, it forms a \emph{closure system}. The associated operator defined by $\phi_{\Sigma}(A) = \bigcap \{ S \in \mathcal{S}_{\Sigma} \mid A \subseteq S \}$ satisfies the axioms of a closure operator (extensivity, monotonicity, and idempotence \cite{fca-book}).
\qed
\end{proof}

\textbf{Structural Constraint Saturation.} 
MLB constraints inherently define a hierarchical (ultrametric) relationship, where $(d_x,d_y,d_z)$ implies that $d_x$ and $d_y$ are more closely related to each other than to $d_z$. However, standard FCA closure operators operate on flat propositional implications (as formalized in Definition~1, $d_x \land d_z \rightarrow d_y$). To bridge this gap, we must translate the structural transitivity of the MLB hierarchy into a complete set of implication rules before applying $\phi_\Sigma$. 

We perform a structural saturation step on $\Sigma$ by exhaustively applying transitive inference rules. For instance, if two constraints $(d_{x_1}, d_{y_1}, d_{z_1})$ and $(d_{x_2}, d_{y_2}, d_{z_2})$ overlap in their sets of strictly closer elements, we dynamically derive and add the implied MLB constraints (e.g., if $d_{a}, d_{b}$ merge before $d_{c}$, and $d_{b}, d_{c}$ merge before $d_{e}$, it implies $d_{a}, d_{b}$ must merge before $d_{e}$). This deductive closure ensures that all indirect hierarchical dependencies are explicitly encoded as implications, allowing the FCA closure operator to natively process the complete underlying tree structure.

\textbf{The Lattice as the Search Space of Valid Clusters.} 
A hierarchical clustering (i.e., dendrogram) is traditionally defined as a nested sequence of clusters.
In the presence of constraints, we define a valid hierarchy $\mathcal{V}$ as follows:

\begin{definition}[Valid Hierarchy] % V for valid hierarchy, H is already used for hierarchy of clusters
    \label{def:valid_hierarchy}
A collection of subsets $\mathcal{V} \subseteq \mathcal{P}(G)$ is a valid hierarchy with respect to $\Sigma$ if: 
\begin{enumerate}
    \item \label{item:valid_hierarchy_root} 
    $G \in \mathcal{V}$ and $\emptyset \in \mathcal{V}$.

    \item \label{item:valid_hierarchy_nested} 
    For all $A, B \in \mathcal{V}$, either $A \cap B = \emptyset$ or $A \subseteq B$ or $B \subseteq A$ (Nested property).
    
    \item \label{item:valid_hierarchy_consistency} 
    Every $A \in \mathcal{V}$ is $\Sigma$-consistent.

    \item \label{item:valid_hierarchy_counterexample}
    For every $(d_x, d_y, d_z) \in \Sigma$, there exists $C \in \mathcal{V}$ such that $d_x, d_y \in C$ and $d_z \notin C$.
\end{enumerate}
\end{definition}

Condition~\ref{item:valid_hierarchy_consistency} and \ref{item:valid_hierarchy_counterexample} of \eqref{def:valid_hierarchy} corresponds to Definitions~\eqref{def:mlb_implication} and \eqref{def:mlb_counterexample}, respectively. 
While the closure operator $\phi_\Sigma$ strictly enforces the implication part (Condition~\ref{item:valid_hierarchy_consistency}), it does not natively enforce the strict structural separation required by Condition~\ref{item:valid_hierarchy_counterexample}. However, we claim that the resulting concept lattice provides the exact \emph{search space} containing all admissible building blocks (clusters) for any valid hierarchy.

\begin{theorem}
Let $\mathfrak{B}_\Sigma$ be the lattice of all closed sets under $\phi_\Sigma$. For any valid hierarchy $\mathcal{V}$, it holds that $\mathcal{V} \subseteq \mathfrak{B}_\Sigma$.
\end{theorem}

\begin{proof}
Let $C \in \mathcal{V}$ be an arbitrary cluster in a valid hierarchy. By the Condition~\ref{item:valid_hierarchy_consistency} of \eqref{def:valid_hierarchy}, $C$ must be $\Sigma$-consistent. By the definition of our closure system $\mathcal{S}_\Sigma$, a set is closed if and only if it is $\Sigma$-consistent. Therefore, $C \in \mathcal{S}_\Sigma$. Since $\mathfrak{B}_\Sigma$ is the lattice formed by all elements of $\mathcal{S}_\Sigma$, it follows that $C \in \mathfrak{B}_\Sigma$.
\end{proof}

This theorem demonstrates that while $\mathfrak{B}_\Sigma$ itself contains overlapping concepts, it encompasses every valid cluster of every possible hierarchy that respects the constraints.
This implies that the concept lattice $\underline{\mathfrak{B}}(\mathbb{K}_{\Sigma})$ constitutes the exact search space for constrained hierarchical clustering. A particular dendrogram can be derived by selecting a sub-order that satisfies Condition~\ref{item:valid_hierarchy_nested} of \eqref{def:valid_hierarchy}, % tree property (nestedness), 
for instance by prioritizing concepts with higher stability or topic coherence.